% =====================================================================
%  1.4 Rasjonale uttrykk — ekstra utfordring
% =====================================================================

\begin{infoboks}{Velg metode selv}
Arket bygger på 1.4 Rasjonale uttrykk. Vis regning og begrunnelse. Bruk eksakte svar.
Skriv vilkår. Kontroller svaret ved å sette inn et tall.
\end{infoboks}

\nivaahode{2}

\oppg{U1}{2}{\uh}
Forenkle. Skriv vilkår.
\[ \frac{x^2-4}{x^2+4x+4} : \frac{x-2}{x^2+2x} \]

\oppg{U2}{2}{\uh}
Skriv som én brøk. Hva forteller svaret om uttrykket?
\[ \frac{1}{x-1}-\frac{1}{x+1}-\frac{2}{x^2-1} \]

\oppg{U3}{2}{\uh}
Forenkle brøken med brøk i teller og nevner. Skriv vilkår.
\[ \dfrac{1+\dfrac{1}{x}}{1-\dfrac{1}{x^2}} \]

\nivaahode{3}

\oppg{U4}{3}{\uh}
Se på uttrykket $\dfrac{x^2+1}{x-1}$.
\begin{tasks}(1)
  \task Vis at $\dfrac{x^2+1}{x-1} = x+1+\dfrac{2}{x-1}$ ved å sette høyre side på felles brøkstrek.
  \task Hva sier omskrivingen om uttrykket når $x$ er et stort tall?
\end{tasks}

\oppg{U5}{3}{\uh}
\begin{tasks}(1)
  \task Vis at $\dfrac{1}{n}-\dfrac{1}{n+1} = \dfrac{1}{n(n+1)}$.
  \task Bruk a) til å finne $\dfrac{1}{1\cdot 2}+\dfrac{1}{2\cdot 3}+\ldots+\dfrac{1}{99\cdot 100}$ eksakt.
\end{tasks}

\oppg{U6}{3}{\prog}
Se på programmet.
\begin{kode}
from fractions import Fraction

sum_ = Fraction(0)
for n in range(1, 11):
    sum_ = sum_ + Fraction(1, n * (n + 1))
    print(n, sum_)
\end{kode}
\begin{tasks}(1)
  \task Hvilket mønster ser du i utskriften?
  \task Forklar hvorfor programmet ikke beviser mønsteret, men U5 a) gjør det.
  \task Endre programmet slik at det summerer $\dfrac{1}{1\cdot 3}+\dfrac{1}{3\cdot 5}+\ldots$ Finn mønsteret og bevis det som i U5.
\end{tasks}

\oppg{U7}{3}{\uh}
For hvilke hele tall $x$ er $\dfrac{x+5}{x-1}$ et helt tall?
\hint{Se på U4.}

\begin{oppsummering}
Faktoriser først, og skriv vilkår før du forkorter. Bruk parentes når det står minus foran en brøk.
Å skrive et uttrykk på en annen form kan vise hva uttrykket egentlig gjør.
\end{oppsummering}
